\documentclass[12pt]{article}
\usepackage[utf8]{inputenc}
\usepackage[T1]{fontenc}
\usepackage{hyperref}
\title{Axi Language Resources for K-12 Curriculum Development}
\author{Compiled by Claude Sonnet 5.5 (claude-sonnet-5-5) for Andr\'as Kornai}
\date{}
\begin{document}
\maketitle

\section{Introduction}
Axi (ISO 639-3 yix; Glottocode axiy1235) is a Loloish language of Yunnan, China, with about 100,000 speakers. It is distinct from Sani, a separate Yi language. The searches found a reference grammar and general academic literature on Yi bilingual education, but no Axi-specific teaching materials.

\section{Primary Educational Resources}
No Axi-specific primers or textbooks were found. A University of Tokyo library page lists early Yi-script materials, including a Yi Script Literacy Textbook. I did not verify which Yi language it serves: \url{https://u-parl.lib.u-tokyo.ac.jp/?p=3392}

\section{Linguistic and Reference Materials}
\begin{itemize}
\item Wikipedia, Axi language: \url{https://en.wikipedia.org/wiki/Axi_language}
\item Wu Zili (2014), \emph{Axi Yiyu gaiyao} (a grammar sketch of Axi), Yunnan Minzu Chubanshe; Glottolog reference entry: \url{https://glottolog.org/resource/reference/id/699775}
\item Wikipedia, Sani language (a different language, listed to avoid confusion): \url{https://en.wikipedia.org/wiki/Sani_language}
\end{itemize}

\section{Religious and Community Resources}
None found.

\section{Digital and Audio Resources}
None found.

\section{Curriculum Development Context}
The language is taught as a subject in some schools according to the reference description, but I found no document describing those programmes. The most relevant background is Zhang and Tsung (2019, Springer), \emph{Bilingual Education and Minority Language Maintenance in China: The Role of Schools in Saving the Yi Language}, which includes case studies of Sani Yi in Yunnan. The searches returned a bookseller page for it, which I did not confirm is that book's page: \url{https://www.hugendubel.de/de/buch_gebunden/-34390859-produkt-details.html}

A paper on Yi-language education, \emph{The Enlightenment of Huayin Qimeng on the Use of Yi Language and Yi Language Education}, also appeared: \url{https://hrmars.com/ijarbss/article/view/26425/The-Enlightenment-of-Huayin-Qimeng-on-the-Use-of-Yi-Language-and-Yi-Language-Education}
\end{document}